\documentclass[../../../include/open-logic-section]{subfiles}

\begin{document}

\olfileid{sth}{cardinals}{zfc}
\olsection{$\ZFC$ تک پہنچ}

حُسنِ ترتیب دا مسلّمہ پا کے اسیں نظریے دے آخری پڑاؤ تک پہنچ گئے آں۔ ہُن~$\ZFC$ لئی لوڑیندے سارے مسلّمے ساڈے کول نیں۔ تفصیل ایہ اے:

\begin{defn}
نظریے $\ZFC$ دے مسلّمے ایہ نیں۔ عنصراں نال تعیّن، اتحاد، جوڑیاں، ذیلی سیٹاں دے سیٹ، لامتناہیت، بنیاد، حُسنِ ترتیب، تے وکھ کرن تے بدل رکھن دے خاکیاں دیاں ساریاں صورتاں۔ دوجے لفظاں وچ، $\ZFC$ وچ~$\ZF$ دے نال حُسنِ ترتیب دا مسلّمہ شامل اے۔
\end{defn}

$\ZFC$ توں مراد \emph{زرمَیلو–فرینکل} دا سیٹ نظریہ اے جیہدے وچ \emph{چون} دا مسلّمہ شامل اے۔ ایہ گل ذرا اوپری لگ سکدی اے، کیوں جے اسیں جیہڑا مسلّمہ پایا اوہدا ناں «حُسنِ ترتیب» سی، «چون» نہیں۔ پر جدوں اسیں اگے {چون} دا مسلّمہ بیان کراں گے، پتا لگے گا کہ حُسنِ ترتیب دا مسلّمہ (~$\ZF$ دے اندر) چون دے مسلّمے دے برابر اے (\olref[choice][woproblem]{thmwochoice} ویکھو)۔ ایس لئی ایہناں وچوں جیہڑا مرضی اپنا «بنیادی» مسلّمہ منّ لئیے۔ تے «$\ZFC$» ناں علمی لکھتاں وچ پوری طرح رواجی اے۔

\end{document}
